% GENERATED FILE. Edit canonical lessons, then run npm run generate:books.
% canonical-source-hash: fnv1a64:8ada3fffcc64d32e
% canonical-lessons: PA-C60-sikkhna, PA-C60-bhullna, PA-C60-chahuna, PA-C60-liauna, PA-C60-bhejna

\chapter{To learn, To forget, To want, To bring, To send}
\label{ch:pa-to-learn-to-forget-to-want-to-bring-to-send}
\hlchaptermodality{60}

\begin{chapteropening}
\textbf{By the end of this chapter:} \emph{I can say to learn, to forget, to want, to bring and to send.}

Complete the last lesson of chapter 60: i can say to learn, to forget, to want, to bring and to send.
\end{chapteropening}

\section[sikkhṇā]{\pa{ਸਿੱਖਣਾ} (sikkhṇā) — to learn}

\label{lesson:PA-C60-sikkhna}



\begin{quote}
Before the new one: say the Punjabi for to sing, then the Punjabi for to dance.
\end{quote}

\subsection*{You'll want to know: \pa{ਸਿੱਖਣਾ}}
\textbf{\pa{ਸਿੱਖਣਾ}} — \emph{sikkhṇā} — \textquotedblleft{}to learn\textquotedblright{}.

\begin{grammarlens}[title={What you've built}]
One more word for this chapter.
\end{grammarlens}

\subsection*{Your turn}
\begin{itemize}
\raggedright
  \item \emph{Say it:} \emph{sikkhṇā}
  \item \emph{Say it:} \emph{sikkhṇā}, once more
  \item \emph{Say it:} say \emph{sikkhṇā}
  \item \emph{Recall:} say the Punjabi for to sing, then the Punjabi for to dance, then say \emph{sikkhṇā} again
\end{itemize}

\begin{tcolorbox}[breakable,colback=teal!4,colframe=teal!35!black,title={Before you move on}]
What does \pa{ਸਿੱਖਣਾ} mean? (\textquotedblleft{}To learn\textquotedblright{}.) Say it once more.
\end{tcolorbox}
\section[bhullṇā]{\pa{ਭੁੱਲਣਾ} (bhullṇā) — to forget}

\label{lesson:PA-C60-bhullna}



\begin{quote}
Before the new one: say the Punjabi for to dance, then the Punjabi for to learn.
\end{quote}

\subsection*{You'll want to know: \pa{ਭੁੱਲਣਾ}}
\textbf{\pa{ਭੁੱਲਣਾ}} — \emph{bhullṇā} — \textquotedblleft{}to forget\textquotedblright{}.

\begin{grammarlens}[title={What you've built}]
One more word for this chapter.
\end{grammarlens}

\subsection*{Your turn}
\begin{itemize}
\raggedright
  \item \emph{Say it:} \emph{bhullṇā}
  \item \emph{Say it:} \emph{bhullṇā}, once more
  \item \emph{Say it:} say \emph{bhullṇā}
  \item \emph{Recall:} say the Punjabi for to dance, then the Punjabi for to learn, then say \emph{bhullṇā} again
\end{itemize}

\begin{tcolorbox}[breakable,colback=teal!4,colframe=teal!35!black,title={Before you move on}]
What does \pa{ਭੁੱਲਣਾ} mean? (\textquotedblleft{}To forget\textquotedblright{}.) Say it once more.
\end{tcolorbox}
\section[chāhuṇā]{\pa{ਚਾਹੁਣਾ} (chāhuṇā) — to want}

\label{lesson:PA-C60-chahuna}



\begin{quote}
Before the new one: say the Punjabi for to learn, then the Punjabi for to forget.
\end{quote}

\subsection*{You'll want to know: \pa{ਚਾਹੁਣਾ}}
\textbf{\pa{ਚਾਹੁਣਾ}} — \emph{chāhuṇā} — \textquotedblleft{}to want\textquotedblright{}.

\begin{grammarlens}[title={What you've built}]
One more word for this chapter.
\end{grammarlens}

\subsection*{Your turn}
\begin{itemize}
\raggedright
  \item \emph{Say it:} \emph{chāhuṇā}
  \item \emph{Say it:} \emph{chāhuṇā}, once more
  \item \emph{Say it:} say \emph{chāhuṇā}
  \item \emph{Recall:} say the Punjabi for to learn, then the Punjabi for to forget, then say \emph{chāhuṇā} again
\end{itemize}

\begin{tcolorbox}[breakable,colback=teal!4,colframe=teal!35!black,title={Before you move on}]
What does \pa{ਚਾਹੁਣਾ} mean? (\textquotedblleft{}To want\textquotedblright{}.) Say it once more.
\end{tcolorbox}
\section[liāuṇā]{\pa{ਲਿਆਉਣਾ} (liāuṇā) — to bring}

\label{lesson:PA-C60-liauna}



\begin{quote}
Before the new one: say the Punjabi for to forget, then the Punjabi for to want.
\end{quote}

\subsection*{You'll want to know: \pa{ਲਿਆਉਣਾ}}
\textbf{\pa{ਲਿਆਉਣਾ}} — \emph{liāuṇā} — \textquotedblleft{}to bring\textquotedblright{}.

\begin{grammarlens}[title={What you've built}]
One more word for this chapter.
\end{grammarlens}

\subsection*{Your turn}
\begin{itemize}
\raggedright
  \item \emph{Say it:} \emph{liāuṇā}
  \item \emph{Say it:} \emph{liāuṇā}, once more
  \item \emph{Say it:} say \emph{liāuṇā}
  \item \emph{Recall:} say the Punjabi for to forget, then the Punjabi for to want, then say \emph{liāuṇā} again
\end{itemize}

\begin{tcolorbox}[breakable,colback=teal!4,colframe=teal!35!black,title={Before you move on}]
What does \pa{ਲਿਆਉਣਾ} mean? (\textquotedblleft{}To bring\textquotedblright{}.) Say it once more.
\end{tcolorbox}
\section[bhejṇā]{\pa{ਭੇਜਣਾ} (bhejṇā) — to send}

\label{lesson:PA-C60-bhejna}



\begin{quote}
Before the new one: say the Punjabi for to want, then the Punjabi for to bring.
\end{quote}

\subsection*{You'll want to know: \pa{ਭੇਜਣਾ}}
\textbf{\pa{ਭੇਜਣਾ}} — \emph{bhejṇā} — \textquotedblleft{}to send\textquotedblright{}.

\begin{grammarlens}[title={What you've built}]
One more word for this chapter.
\end{grammarlens}

\subsection*{Your turn}
\begin{itemize}
\raggedright
  \item \emph{Say it:} \emph{bhejṇā}
  \item \emph{Say it:} \emph{bhejṇā}, once more
  \item \emph{Say it:} say \emph{bhejṇā}
  \item \emph{Recall:} say the Punjabi for to want, then the Punjabi for to bring, then say \emph{bhejṇā} again
\end{itemize}

\begin{tcolorbox}[breakable,colback=teal!4,colframe=teal!35!black,title={Before you move on}]
What does \pa{ਭੇਜਣਾ} mean? (\textquotedblleft{}To send\textquotedblright{}.) Say it once more.
\end{tcolorbox}
